\documentclass[UTF8,zihao=-4]{ctexart}

\usepackage[a4paper,top=2.4cm,bottom=2.4cm,left=2.25cm,right=2.25cm]{geometry}
\usepackage{amsmath,amssymb}
\usepackage{booktabs,longtable,array,tabularx}
\usepackage{listings}
\usepackage{enumitem}
\usepackage{xcolor}
\usepackage{hyperref}
\usepackage{fancyhdr}

\definecolor{inkblue}{RGB}{34,76,112}
\definecolor{softgray}{RGB}{246,248,250}
\lstset{basicstyle=\small\ttfamily,breaklines=true,showstringspaces=false,frame=single,rulecolor=\color{softgray}}
\newcommand{\slot}[1]{\textcolor{inkblue}{\textit{【#1】}}}
\newcommand{\guidehead}[1]{\par\medskip\noindent\textbf{\color{inkblue}#1}\par}
\hypersetup{
  colorlinks=true,
  linkcolor=inkblue,
  urlcolor=inkblue,
  citecolor=inkblue,
  pdfauthor={数学建模学习小组},
  pdftitle={数学建模阶段学习笔记与论文写作模板}
}

\pagestyle{fancy}
\fancyhf{}
\fancyhead[L]{数学建模阶段学习笔记与论文写作模板}
\fancyhead[R]{\thepage}
\setlength{\headheight}{15pt}
\setlength{\parindent}{2em}
\setlength{\parskip}{0.25em}
\setlist{nosep,leftmargin=2.2em}

\title{\bfseries 数学建模阶段学习笔记与论文写作模板}
\author{队长：夏同、组员1：温启韬、组员2：梁文曦}
\date{2026年8月11日}

\begin{document}

\maketitle

\begin{abstract}
这份报告是我们现阶段队内训练的记录，同时也是在队长群内由刘小兰教授提出，丁老师转发的作业。

第一部分，我们依照思维导图和三人分工，整理了十三个算法与模型学习单元，并借助公开资料和编程工具核对代码逻辑，把原理、求解步骤、代码入口和使用边界放在一起，形成阶段学习报告；

第二部分，我们选取2024年A题一等奖论文A053《基于目标优化的板凳龙形态调节》，细看摘要、假设、建模、求解、结果解释和附录怎样前后接续；同时把2020--2025年A、B、C题59篇编号论文的章节、篇幅和判断动作作为交叉依据。A053负责说明一篇论文怎样展开，59篇语料负责约束我们模板的共同骨架、篇幅范围和可选写法。

\end{abstract}

\tableofcontents
\newpage

\section{阶段学习笔记}

\subsection{选题侧重与成员分工}

近六年的A、B、C题考查的能力各有侧重。A题往往要从物理过程落到方程，还得看数值误差；B题的要害大概率在变量、约束和可行方案；C题通常先处理数据口径，再谈预测、评价和决策。我们组分工按任务侧重展开：队长（夏同）主攻机理与连续模型，组员1（温启韬）负责规划和优化，组员2（梁文曦）负责数据与统计学习。我们在学习的过程中尽量争取能用自己的话解释为什么选它，能把公式拆成程序上的一套可以搬运或者模仿的动作，也能说出它在哪些条件下会失效。

但是，实在由于思维导图中的方法远不止这些，我们组这次先完成十三个单元的原理与代码梳理；尚未动手的分支我们后面用真题逐项补齐，或者计划以赛代练。

\subsection{十三个算法与模型学习单元}

\setlength{\tabcolsep}{2.5pt}
\small
\begin{longtable}{p{1.25cm}p{2.55cm}p{6.25cm}p{5.45cm}}
\caption{成员学习内容、代码逻辑与使用边界}\label{tab:learning-units}\\
\toprule
成员 & 算法或模型 & 基本原理、步骤与代码逻辑 & 适用问题、优缺点和限制\\
\midrule
\endfirsthead
\toprule
成员 & 算法或模型 & 基本原理、步骤与代码逻辑 & 适用问题、优缺点和限制\\
\midrule
\endhead
\bottomrule
\endfoot

夏同 & ODE与RK4
& 先从受力、守恒或状态转移关系写出 $\dot{y}=f(t,y)$，再由初值出发，用四个斜率的加权值向前推进。
& 适合状态随时间连续变化的力学、传热、种群和传播过程。RK4在精度与实现难度之间较为均衡，但刚性方程会受稳定域限制；此时应改用刚性求解器，不靠缩小一步掩盖问题。\\

夏同 & PDE与有限差分
& 空间区域、初始条件和边界条件要先落到网格上，再将时间、空间导数换成差分格式。
& 适合规则区域上的热、波和扩散问题。公式与网格位置对应直观，代价是复杂边界难处理，显式格式还须核对稳定条件。\\

夏同 & 有限元模型
& 从PDE的形式入手：划分单元，计算局部矩阵，装配全局稀疏矩阵，“施加”边界后求节点量。
& 适合不规则几何、变系数介质和混合边界。几何适应性强，但网格质量、形函数和边界
装配都容易引入错误；规则区域的入门题先用有限差分作基线。\\

温启韬 & LP与MILP
& 先逐项列出决策变量、线性目标、等式和不等式；遇到启停、选择或分配，才引入0-1变量。
& 适合资源分配、排程、选址和生产决策。线性模型便于解释并能给最优性信息；整数规模大时计算量上升，Big--M须来自实际上界，不能任意指定过大的常数。\\

温启韬 & 粒子群优化
& 每个粒子都保留位置、速度、个体最好解和群体最好解；迭代时包含速度更新、越界处理和重新评价。
& 适合连续、非光滑或难求导的参数搜索，编码简短；它没有连续全局最优证书，结果会受边界、罚项和随机种子影响。若变量只有一维且单调，可改用遍历或二分。\\

温启韬 & 图论与网络优化
& 建模时把对象落成节点、关系落成带权边。非负单源最短路用Dijkstra；容量与费用同时出现时用最小费用流。实战时，先建图，再运行算法，最后把路径或流量映射回现实对象并核对容量。
& 适合运输、路径、匹配和网络资源配置。结构清楚、结果易复算；动态图、负权边和多目标路线要另行处理，最小生成树也不能代替两点间最短路。\\

梁文曦 & 缺失、异常与重复处理
& 拿到附件后，先建数据字典和主键，并把未检出、未知、未采样分开记录；随后才按成因决定删除、插值或KNN。
& 它决定后续样本是否可信，但不能在划分训练集前用全体数据估计均值或尺度。\\

梁文曦 & PCA与K--Means
& 数据标准化后，对协方差矩阵作特征分解并按解释率选取主成分；降维结果再送入K--Means，比较不同簇数的轮廓系数。
& PCA适合线性降维，K--Means适合近似球形簇。两者速度快、图形直观。\\

梁文曦 & 熵权与TOPSIS
& 指标须先正向化和无量纲化。熵权刻画样本差异，TOPSIS据此计算对象到正、负理想解的距离。
代码逐步导出标准矩阵、权重、距离和排序，并对权重作扰动。
& 适合多指标评价与方案排序。计算过程清楚，但“差异大”不等于“业务上更重要”；指标
方向、相关重复和权重敏感性都要说明，不能只给最终名次。\\
\end{longtable}
\normalsize

\subsection{本阶段做得较好的地方}

我们三人的工作成果能够在同一变量体系中完美衔接、相互适配。同时，我们搭建了多组基础对照基线，保障算法验证有据可依。此外，在研读群内上传的论文等资料时，我们不再只梳理目录或者大体框架，而是梳理出论文的完整逻辑流程，并整理成可直接对应、核查代码的实操笔记，落地性更强。

\subsection{目前不足}

仅凭已有的学习成果而不落地实操，并不能确保团队成员都能在比赛规定时间内独立完成代码调试、跑通完整流程。

\section{一等奖论文学习分析}

\subsection{选文与阅读范围}

本次选择2024年A题《基于目标优化的板凳龙形态调节》\cite{A053}。

\subsection{论文结构及各部分写法}

\small
\begin{longtable}{p{2.25cm}p{4.0cm}p{5.1cm}p{4.2cm}}
\caption{A053章节功能与我们理解}\label{tab:a053-structure}\\
\toprule
部分 & 论文在写什么 & 具体写法 & 我们吸收的做法\\
\midrule
\endfirsthead
\toprule
部分 & 论文在写什么 & 具体写法 & 我们吸收的做法\\
\midrule
\endhead
\bottomrule
\endfoot
摘要（p.1）
& 依次回答运动状态、碰撞、最小螺距、调头曲线和速度上限。
& 每个分问都给对象、模型动作和交付值，如第一次碰撞时刻、螺距和龙头速度；算法名放在它所求的量之后。
& 我们写摘要时也按分问落笔，但每问至少留下一个能作判断的结果；顺序词能省则省，不把摘要写成目录。\\

问题分析（p.3--4）
& 把板凳龙的整体运动拆成相邻把手位置递推，再说明各问增加的约束。
& 问题四被拆成“求较短调头曲线”和“求调头过程中的运动状态”两个子任务，前者输出路径，后者沿用位置递推。
& 我们把输入输出关系画清后，再决定方法；沿用前问时，直接列明哪些量继承、哪些量新加。\\

模型假设（p.2）
& 限定人为扰动、板凳刚体和开始调头的时机。
& 假设会进入方程或事件判定，没有写与计算无关的常识句。
& 每条假设都要回答“它使哪一项可以怎样计算”，并在评价中说明放宽后要改哪个模块。\\

符号说明（p.3）
& 给出螺距、极角、位置、速度、圆弧半径等贯穿五问的量。
& 统一公共符号，分问新增量在首次使用处定义，减少同一字母承担两个含义。
& 符号表只放多次使用的量；仅在一处使用的中间量随公式解释。\\

模型建立（p.6--21、p.23--36）
& 写螺线弧长、定长递推、矩形碰撞、相切圆弧和多段路径状态。
& 公式按“已知把手--定长方程--候选根--运动方向筛选--下一把手”推进；碰撞从几何对象落到矩形投影区间。
& 公式尽量按程序实际执行的次序出现。多根问题交代筛选条件，分段路径则先列全状态及切换条件。\\

模型求解（p.8--16、p.22--31、p.34--41）
& 用二分、时间离散、变步长遍历、状态递推和粒子群得到五问结果。
& 螺距先在 $[45,46]\,\mathrm{cm}$ 内定位，再以 $0.01\,\mathrm{cm}$ 步长细查；碰撞检测在第一次触发时停止。
& 粗查负责圈定“边界在哪一段”，细查再回答“边界精确到多少”；随机算法另报种子和重复次数。\\

结果及分析（p.22、p.27--28、p.32--33、p.38、p.41）
& 报告碰撞时刻、螺距、圆弧半径、曲线长度和速度上限，并解释局部异常。
& 参数微变后若第一次碰撞对象改变，论文从局部几何解释输出跳变；12--13秒的速度突增则用相同时间内的路程差解释。
& 表后的正文不复述数字，而要回答“为何变、哪条约束起作用、结论只到什么范围”。\\

检验与边界
& 检查首次碰撞对象、全过程可行性和现实动作限制。
& p.17用反设和逐节前推说明最早碰撞对象；p.28用“终点可行但中途碰撞”的方案否定只查终点；p.32从龙身不能倒退推出半径下界。
& 检验要对准一个具体疑问。回代、反例和物理边界各有侧重。\\

附录（p.43--62）
& 给出逐把手循环、二分、分段路径、碰撞和粒子群代码。
& 正文四种运动状态与代码状态码1--4相互对应，表格输出也由固定变量名生成。
& 附录须留下复现正文结果所需的入口、参数和输出说明；正文只讲关键流程，不整段粘贴程序。\\
\end{longtable}
\normalsize

\subsection{五个分问的思路怎样连起来}

\begin{enumerate}
  \item 问题一先确定龙头在阿基米德螺线上的位置，相邻把手随后按距离不变的关系逐节递推；
  速度由短时间内的弧长变化得到。这一问最后既留下数值，也留下后四问共用的位置、速度计算器。
  \item 问题二把每块板凳写成矩形，用分离轴判断重叠。这里没有一上来就缩小检查范围，
  而是先用反证法说明第一次碰撞对象，再让程序在第一次触发处停止。
  \item 到了问题三，“能进入调头区域”被落实为全过程均不碰撞。螺距从粗区间缩到细步长，
  终点可行却中途碰撞的候选被舍去，所以可行性不是只看最后一帧。
  \item 问题四先由相切关系写出两段圆弧，再用“后续构件不能倒退”给大圆弧半径下界；
  确定路径后，位置递推按盘入、两段圆弧和盘出四种状态继续运行。
  \item 问题五不另建运动模型，只把问题四所得的各把手速度改写成对于上限的约束；粒子群在这里只负责
  搜索龙头速度。
\end{enumerate}

\subsection{最值得模仿的段落组织}

\paragraph{从局部关系写到全局运动。}
板凳龙的构件很多，真正反复出现的关系却只有相邻把手间的固定距离。一个把手的位置求出后，下一把手的位置转化为一组定长方程；得到的根还要经构件次序、运动方向的筛选。

\paragraph{先证明检查范围不会漏，再缩小程序。}
讨论第一次碰撞时，论文先假设第 $i$ 个构件更早碰撞，再借助轨迹的先后关系逐节向前追，直到同龙头运动相矛盾。这段证明承担的职责是“不漏掉更早事件”。有了这个结论，程序才有理由只查与龙头有关的候选对象。今后若要剪枝，我们也应先说明结构理由，再谈省了多少计算。

\paragraph{用反例判断并简化方案。}
螺距约为 $42\,\mathrm{cm}$ 的方案在终点附近看似可行，但在完整运动过程中会提前碰撞。这个反例把“只查终点”的漏洞钉在了运动途中：候选方案到达终点前便已碰撞。当速度在12--13秒突然抬升时，论文比较同一时间内相邻节点走过的路程，并沿小圆弧附近的路径形状查找原因。它的次序是圈出异常时段，定位局部位置，按同一口径比较，最后才下速度结论。


\section{本队论文模板（可填写版）}

本节不是“写作建议清单”，而是一份可以复制到正式论文中逐项填写的稿件蓝图。它由工作区内 2020--2025 年 A、B、C 题 59 篇编号论文、共 2892 页的结构核对生成\cite{corpus59}；A053 只用于上一节的单篇精读。模板的目录顺序固定，判断的起笔顺序不固定：真实论文可以从一张图、一个边界、一次失败试算或前一问留下的接口开始，只要把取舍的依据和后果交代清楚。蓝色的\slot{填空位}是训练阶段的占位符，提交前须全部替换或删除。

\subsection{先按竞赛规章搭好提交版式}

以下内容逐条核对自工作区新增的三份全国组委会文件：《全国大学生数学建模竞赛论文格式规范（2026年修订稿）》《全国大学生数学建模竞赛参赛规则（2026年修订稿）》和《全国大学生数学建模竞赛人工智能工具使用规定（2026年试行）》\cite{format2026,rules2026,ai2026}。本学习报告为了呈现作业内容保留目录和队员姓名；真正参赛时必须另建提交版，不能把本报告原样交赛。

\begin{longtable}{p{2.7cm}p{5.7cm}p{6.2cm}}
\caption{2026 年论文与支撑材料硬性要求}\label{tab:official-format-rules}\\
\toprule
对象 & 纸质版/正文要求 & 电子版/支撑材料要求\\
\midrule
\endfirsthead
\toprule
对象 & 纸质版/正文要求 & 电子版/支撑材料要求\\
\midrule
\endhead
纸张和装订 & 白色 A4 纸，单面或双面均可；四边页边距至少 $2.5\,\mathrm{cm}$，左侧装订。 & 电子论文必须与纸质版内容及格式（含附录）一致。\\
首页顺序 & 第 1 页承诺书，第 2 页编号专用页，第 3 页摘要专用页。 & 电子论文不放承诺书和编号专用页，第一页直接是摘要专用页。\\
摘要与页码 & 摘要含标题、关键词，无需英文，原则上不超过 1 页；从摘要页起以阿拉伯数字从 1 连续编号，页码居页脚中部。 & 参赛论文为一个单独 PDF 或 Word 文件，建议 PDF，文件不超过 20 MB，不能再压缩后上传。\\
正文与附录 & 正文从纸质版第 4 页开始，不设目录，不超过 30 页；正文后为附录，页数不限，须打印并与正文装订。 & 支撑材料另交一个 RAR 或 ZIP 文件，不超过 20 MB。\\
附录和程序 & 附录列出支撑材料文件，并纳入全部完整、可运行的源程序（含 Excel、SPSS 等软件的交互命令）；确实未用程序时写明“本论文没有用到程序”。 & 支撑材料至少包括全部可运行源程序、自主查阅数据（赛题原始数据除外）和较大篇幅中间图表；确实没有时在附录注明“本论文没有支撑材料”。\\
匿名与引用 & 摘要、正文和附录不得出现参赛者身份、学校或赛区信息；他人及公开资料须在正文标注并列入参考文献。 & 支撑材料同样不得出现身份、学校或赛区信息，也不得放入承诺书和编号专用页。\\
\bottomrule
\end{longtable}

竞赛期间必须由本队队员独立完成赛题。指导教师不能解释赛题、提供选题或解题建议、提供资料、修改论文或给出修改意见；队员也不能同队外任何人讨论赛题，不能在社交平台、论坛、代码平台或直播间浏览、发布、讨论当届赛题内容。AI 工具可以作为辅助，但作品的原创性、真实性和准确性仍由参赛队负责，并须遵守专门披露规定。

\subsubsection{AI 工具披露位置与材料}

2026 年试行规定要求参赛队公开、如实说明 AI 工具的使用情况，并逐项人工审查和核实相关内容。声明不是附录，也不放在参考文献之后，而是作为独立部分排在参考文献之前；两种情况只能选择其一。

\begin{longtable}{p{2.7cm}p{5.7cm}p{6.2cm}}
\caption{AI 工具使用披露清单}\label{tab:official-ai-rules}\\
\toprule
实际情况 & 论文内处理 & 支撑材料与人工责任\\
\midrule
\endfirsthead
\toprule
实际情况 & 论文内处理 & 支撑材料与人工责任\\
\midrule
\endhead
未使用任何 AI 工具 & 在参考文献之前设置“AI 工具使用声明”，使用规定给出的未使用声明。 & 不提交虚构记录；声明必须与竞赛期间的实际过程一致。\\
使用 AI 工具 & 在参考文献之前设置“AI 工具使用声明”，简要写明用途，并指向支撑材料。 & 支撑材料必须包含独立 PDF，文件名为\texttt{AI 工具使用详情.pdf}；写明工具名称及版本或型号、使用目的和环节、主要提示方式与使用过程、采纳情况以及人工修改与核验。语言润色可不写最后一项。\\
核心建模与分析 & 核心建模与分析必须由参赛队主导；AI 参与完成的内容须逐项人工审查与核实。 & 故意隐瞒、虚假声明，或把未经必要人工审查和核实的 AI 生成内容直接作为核心成果提交，将被取消评奖资格。\\
\bottomrule
\end{longtable}

正式论文在参考文献前保留下面这一段，按实际情况删去另一项：
\begin{quote}\small
\textbf{AI 工具使用声明}

未使用时：本参赛队在竞赛过程中未使用任何 AI 工具。

使用时：本参赛队在竞赛过程中使用了 AI 工具，主要用于\slot{简要用途，如语言润色、代码调试等}，详细使用情况见支撑材料。
\end{quote}
规定没有要求把 AI 工具另列为参考文献，也没有要求在正文中逐句标出 AI 辅助痕迹；但是支撑材料中的用途、过程和人工核验必须能够与论文及代码相互对应。该规定自 2026 年 9 月 1 日起试行，此前规定与之不一致时以本规定为准。

\subsection{全量语料提取出的目录与篇幅基准}

下面的页数是页级标题跨度的中位数及四分位范围，不是硬性配额。结果分析和附录本来就随题目、图表和代码长度变化；模型检验、灵敏度和改进方案在语料中的检出率较低，只能作为“有证据才写”的功能要求，不能倒推每篇都必须有同样的小节。

\begin{longtable}{p{2.3cm}p{1.5cm}p{3.0cm}p{8.0cm}}
\caption{59 篇 A/B/C 论文的章节跨度与交付物}\label{tab:template-corpus-length}\\
\toprule
章节 & 检出数 & 页级跨度中位数 [Q1,Q3] & 正式稿中必须留下的可核对内容\\
\midrule
\endfirsthead
\toprule
章节 & 检出数 & 页级跨度中位数 [Q1,Q3] & 正式稿中必须留下的可核对内容\\
\midrule
\endhead
摘要与关键词 & 58/59 & 1 [1,1] & 分问任务、方法动作、关键数值、检验口径、结论和 3--5 个关键词。\\
问题重述 & 57/59 & 1 [1,1] & 研究对象、输入、输出、硬约束和分问之间的接口。\\
问题分析 & 57/59 & 1 [1,1] & 从题面、数据或物理边界出发的取舍依据，以及为何保留或放弃某条路线。\\
模型假设 & 48/59 & 1 [1,1] & 每条假设的对象、条件、计算后果和适用边界。\\
符号说明 & 58/59 & 1 [1,3.75] & 变量、参数、集合、下标、单位；只放跨公式复用的量。\\
模型建立 & 58/59 & 2 [1,6] & 定义、方程/目标函数、约束、边界条件和与题意的映射。\\
模型求解 & 55/59 & 2 [1,6] & 初始化、更新、可行性处理、搜索范围、停止条件、输出和代码接口。\\
结果分析 & 47/59 & 7 [1,12] & 读数、方向、数量级、机制、约束状态、现实含义和结果边界。\\
模型检验 & 22/59 & 1.5 [1,7] & 检验对象、独立证据、指标/阈值、结果和裁决；没有检验就不能写“验证有效”。\\
灵敏度分析 & 25/59 & 1 [1,2] & 基准、扰动来源、范围、固定量、输出变化和决策是否切换。\\
模型评价 & 51/59 & 1 [1,1] & 有公式或数据支撑的优点，以及受对象、数据和计算预算限制的缺点。\\
改进方案 & 19/59 & 1 [1,1] & 新增数据/关系、修改层、预计改善量和仍需检验的边界。\\
参考文献 & 57/59 & 1 [1,11] & 正文真正使用过的题面、方法、数据和软件来源。\\
附录 & 48/59 & 16 [6,28] & 代码入口、参数、数据字典、运行顺序、补充推导、全量结果和输出文件。\\
\bottomrule
\end{longtable}

\subsection{正式稿标题树与分问接口}

按下列顺序建立目录；“问题一”等只替换为题目中的实际编号和任务名，不要把算法名称当作标题。每个分问都要在上一问的输出处写清接口：

\begin{enumerate}
  \item 摘要与关键词；
  \item 问题重述；
  \item 问题分析；
  \item 模型假设；
  \item 符号说明；
  \item 问题一：\slot{任务名}；问题二：\slot{任务名}；……每问均包含“模型建立”与“模型求解”；
  \item 结果分析与模型检验；
  \item 灵敏度/稳健性分析（若题目或结果支持）；
  \item 模型评价与改进；
  \item AI 工具使用声明：按实际情况选择官方规定中的一种表述；
  \item 参考文献；
  \item 附录：代码、数据、补充结果和运行说明。
\end{enumerate}

\begin{longtable}{p{1.2cm}p{3.0cm}p{3.0cm}p{3.0cm}p{4.2cm}}
\caption{分问任务矩阵：先填接口，再写正文}\label{tab:template-task-matrix}\\
\toprule
分问 & 输入/证据 & 要求输出 & 硬约束或口径 & 交给下一问的接口\\
\midrule
\endfirsthead
\toprule
分问 & 输入/证据 & 要求输出 & 硬约束或口径 & 交给下一问的接口\\
\midrule
\endhead
问题一 & \slot{附件字段、初始状态、题面量} & \slot{数值/分类/曲线/方案} & \slot{单位、范围、边界} & \slot{变量、参数、筛选结果}\\
问题二 & \slot{沿用问题一的哪些输出} & \slot{……} & \slot{新增约束} & \slot{……}\\
问题三 & \slot{……} & \slot{……} & \slot{……} & \slot{最终决策或评价口径}\\
问题四（如有） & \slot{……} & \slot{……} & \slot{……} & \slot{……}\\
\bottomrule
\end{longtable}

\subsection{摘要与关键词：一页内交付答案}

\guidehead{本节填什么}写清对象和条件、每个分问采用的实际模型动作、最能区分方案的数值或关系、检验口径以及结论边界。不要把摘要写成方法名目录，也不要把尚未在正文或代码中产生的数字填进去。

\guidehead{可直接改写的正文骨架}
\begin{quote}\small
围绕\slot{对象}在\slot{条件/约束}下的\slot{现实任务}，本文先由\slot{题面结构、附件读数或物理边界}确定\slot{变量域/初始状态/数据口径}。对于问题一，\slot{由什么证据触发了哪一个模型动作}，得到\slot{关键结果、单位和范围}；对于问题二，\slot{沿用或修正的问题一接口}，在\slot{新增约束/不确定性}下求得\slot{关键结果}；对于问题三，\slot{评价、预测或决策动作}给出\slot{最终方案/排序/预测}。通过\slot{回代、残差、交叉验证、基线、反例、物理边界或共同环境复评}检查后，\slot{最有判别力的结论}在\slot{适用范围}内成立。\\
关键词：\slot{任务词}；\slot{模型词}；\slot{求解词}；\slot{结果词}
\end{quote}

\guidehead{摘要验收}每个分问至少有一个可复核输出；数字能追到正文表/代码；单位前后一致；检验方法没有被夸大；“最优、稳定、准确”等词后面紧跟比较对象、指标和范围。

\subsection{问题重述：把题面改成任务矩阵}

\guidehead{可直接改写的正文骨架}
题目研究的是\slot{对象及其现实背景}。已知\slot{输入数据、初始状态、时间范围和单位}，同时受到\slot{硬约束、边界条件和可行域}限制。问题一要求\slot{给定什么、求什么}；问题二在\slot{问题一输出/新增数据}基础上要求\slot{……}；问题三进一步比较或决定\slot{……}。本文将\slot{不能混用的口径，如预测值与观测值、局部指标与全局指标}分别记账。

\guidehead{本节不写}不在重述中提前宣布“采用某算法”，不改变题面中的对象范围、时间窗、单位和“必须/可以”等语气，不把附件中没有的字段补成已知量。

\subsection{问题分析：把选择写成有来由的判断}

论文需要让读者看见“为什么这样做”，但没有必要保留未经整理的试算流水账。正文只留下可由题面、数据、公式、图表或运行记录复核的判断，也不强行把每段写成“困难--基线--不足--候选--选择”的固定顺序；可以从一个异常读数、一个不满足的边界、一次失败试算、一个前问接口或一条守恒关系起笔。

\begin{longtable}{p{2.0cm}p{3.5cm}p{3.4cm}p{3.3cm}p{3.4cm}}
\caption{问题分析台账：把真实的取舍留在稿件中}\label{tab:template-analysis-ledger}\\
\toprule
分问 & 观察到的事实/限制 & 尚未解决的量或矛盾 & 采取的动作及其改变 & 交给模型的边界/接口\\
\midrule
\endfirsthead
\toprule
分问 & 观察到的事实/限制 & 尚未解决的量或矛盾 & 采取的动作及其改变 & 交给模型的边界/接口\\
\midrule
\endhead
问题一 & \slot{图、表、物理条件或失败试算} & \slot{要解释/求解的量} & \slot{保留、舍弃、改写或分段} & \slot{变量、范围、初值}\\
问题二 & \slot{问题一的残差/新约束} & \slot{……} & \slot{……} & \slot{……}\\
问题三 & \slot{决策冲突或不确定性} & \slot{……} & \slot{……} & \slot{……}\\
\bottomrule
\end{longtable}

\guidehead{可直接改写的连接句}“附件中\slot{某字段/某段轨迹}先呈现出\slot{具体形态}，这使得\slot{某种简化/某个候选}不能直接承担\slot{任务}。因此先保留\slot{必须保留的关系}，把\slot{新增自由度/误差来源}留到\slot{后续步骤}处理。这样得到的输出\slot{如何进入下一问}，而不会把\slot{未观测量/事后汇总量}误当成决策时可用信息。”

\subsection{模型假设：每条假设都要产生计算后果}

\begin{longtable}{p{3.2cm}p{3.1cm}p{4.0cm}p{3.6cm}}
\caption{模型假设台账}\label{tab:template-assumption-ledger}\\
\toprule
假设（对象+条件） & 题面/数据依据 & 在公式或程序中的后果 & 检查与失效边界\\
\midrule
\endfirsthead
\toprule
假设（对象+条件） & 题面/数据依据 & 在公式或程序中的后果 & 检查与失效边界\\
\midrule
\endhead
\slot{忽略某项影响} & \slot{题面允许/量级比较} & \slot{删除哪一项或固定哪个量} & \slot{在何种范围失效}\\
\slot{连续/刚体/独立等条件} & \slot{观测或定义} & \slot{改变哪种状态转移或误差} & \slot{用何种回代/扰动检查}\\
\bottomrule
\end{longtable}

正文句式：“在\slot{条件}下，\slot{对象}的\slot{因素}相对于\slot{基准}可忽略，故在模型中\slot{具体处理}。若\slot{边界条件}不再满足，该处理只适用于\slot{范围}，需要在\slot{改进/敏感性}中重新计算。”

\subsection{符号说明：只收录会反复使用的量}

\begin{longtable}{p{2.2cm}p{7.0cm}p{2.2cm}p{3.3cm}}
\caption{符号说明可填写表}\label{tab:template-symbol-ledger}\\
\toprule
符号 & 含义（对象和方向） & 单位/取值域 & 首次出现位置\\
\midrule
\endfirsthead
\toprule
符号 & 含义（对象和方向） & 单位/取值域 & 首次出现位置\\
\midrule
\endhead
$x_i$ & \slot{第 $i$ 个对象的决策/观测量} & \slot{单位或域} & 式(\slot{})\\
$\theta$ & \slot{待估参数/情景参数} & \slot{单位或范围} & 式(\slot{})\\
$\mathcal{I}$ & \slot{对象集合/时间集合} & \slot{集合定义} & 式(\slot{})\\
\bottomrule
\end{longtable}

\subsection{分问题模型建立与求解：每问一张模型卡}

先复制下面的模型卡，再为每一问填写。模型名称不能替代模型内容；若用了现成求解器，仍须说明它接收什么、保持什么、何时停止和输出什么。

\begin{longtable}{p{1.8cm}p{2.4cm}p{2.7cm}p{2.7cm}p{2.8cm}p{2.1cm}}
\caption{分问模型卡}\label{tab:template-model-card}\\
\toprule
分问 & 变量/状态 & 目标或方程 & 约束/边界 & 求解动作 & 输出及核查\\
\midrule
\endfirsthead
\toprule
分问 & 变量/状态 & 目标或方程 & 约束/边界 & 求解动作 & 输出及核查\\
\midrule
\endhead
问题一 & \slot{} & \slot{} & \slot{} & \slot{初始化--更新--停止} & \slot{}\\
问题二 & \slot{} & \slot{} & \slot{} & \slot{} & \slot{}\\
问题三 & \slot{} & \slot{} & \slot{} & \slot{} & \slot{}\\
\bottomrule
\end{longtable}

\subsubsection*{问题\slot{编号}：\slot{任务名}}
\paragraph{输入与变量。}输入为\slot{附件/前问输出/外部给定}。令\slot{变量}表示\slot{含义}，其取值域为\slot{范围、符号和单位}。\slot{不参与决策但需固定的量}保持\slot{固定规则}。

\paragraph{模型建立。}由\slot{守恒、几何关系、统计定义、业务规则或题面约束}得到
\begin{equation}
  \slot{目标函数、状态方程或评价式}.
\end{equation}
其中\slot{左侧/右侧各项}分别表示\slot{可解释含义}。可行域由\slot{约束1}、\slot{约束2}和\slot{边界条件}给出；\slot{某约束}直接对应题面中的\slot{原词}。

\paragraph{求解与实现。}先\slot{清洗/初始化/粗略定位}，随后\slot{二分、递推、枚举、线性规划、随机搜索、拟合或仿真}。每次更新保留\slot{状态/可行标志/候选解}，越界时\slot{处理规则}；在\slot{停止条件}满足时输出\slot{文件、表格、曲线或参数}。代码入口为\texttt{\slot{file.py/main.m}}，关键函数为\texttt{\slot{function}}，随机过程使用\slot{种子、重复次数和预算}。

\paragraph{模型边界。}该模型回答的是\slot{对象、时间窗和数据条件}，不能外推到\slot{未观测对象/不同约束/更长时间}。若\slot{观察到的残差或失败状态}出现，则保留\slot{基线/旧模型}并只修改\slot{具体关系或参数}，再按同一口径复算。

\subsubsection{常用公式组件（按题意选用，不要整段照搬）}

\paragraph{连续过程。}状态量、输入和观测量可先写成
\begin{equation}
\dot{\boldsymbol{x}}(t)=\boldsymbol{f}(\boldsymbol{x}(t),\boldsymbol{u}(t),t),\qquad
\boldsymbol{y}(t)=\boldsymbol{g}(\boldsymbol{x}(t),\boldsymbol{u}(t),t),\qquad
\boldsymbol{x}(0)=\boldsymbol{x}_0.
\end{equation}
随后补上\slot{初值、边界、守恒/受力关系}，并说明采用\slot{RK4/FDM/FEM/其他}时步长或网格怎样进入误差控制。公式后必须把\slot{每一项的物理含义}翻译成文字。

\paragraph{规划与优化。}若决策变量为\slot{$x$}，目标和可行域可整理为
\begin{equation}
\begin{aligned}
\min_{x}\quad & F(x)\\
\text{s.t.}\quad & g_i(x)\leq 0,\quad i=1,\ldots,m,\\
& h_j(x)=0,\quad j=1,\ldots,p,\\
& x\in\mathcal{X}.
\end{aligned}
\end{equation}
若有启停、分配或选择，单独说明\slot{0--1变量与其联动约束}；若目标有优先级，写出\slot{字典序、分层优化或权重来源}，不要只给一个未经解释的综合分数。

\paragraph{线性/整数模型。}线性情形可以使用
\begin{equation}
\min\;\boldsymbol{c}^{\mathsf T}\boldsymbol{x}\quad
\text{s.t.}\quad A\boldsymbol{x}\leq\boldsymbol{b},\quad
\boldsymbol{x}\geq0,\quad x_k\in\mathbb{Z}\ \text{(若题意要求)}.
\end{equation}
正文说明\slot{$A,b,c$各列如何由题面生成}、求解器返回的状态以及最终解的约束残差。

\paragraph{统计/评价模型。}预测、分类或评价时，将\slot{原字段--清洗--特征--标签--划分}写成一条数据谱系，再给出损失或指标。例如
\begin{equation}
\mathrm{MAE}=\frac{1}{n}\sum_{i=1}^{n}|y_i-\hat y_i|,\qquad
\mathrm{RMSE}=\sqrt{\frac{1}{n}\sum_{i=1}^{n}(y_i-\hat y_i)^2}.
\end{equation}
指标方向、样本角色和业务解释必须跟代码中实际计算的对象一致。

\subsection{A/B/C 题型专用插槽}

\begin{longtable}{p{1.6cm}p{4.4cm}p{4.5cm}p{4.3cm}}
\caption{三类题型必须补齐的内容}\label{tab:template-abc-slots}\\
\toprule
题型 & 建模时必须写清 & 求解时必须写清 & 结果与检验必须写清\\
\midrule
\endfirsthead
\toprule
题型 & 建模时必须写清 & 求解时必须写清 & 结果与检验必须写清\\
\midrule
\endhead
A：机理/连续 & 状态、坐标、初值、边界、守恒或受力关系；离散格式与连续模型分开。 & 步长/网格、稳定条件、初值推进、事件触发、收敛或误差界；多根按物理顺序筛选。 & 回代、守恒、边界和量纲；异常时段回到局部轨迹或同时间路程，不能只报一条曲线。\\
B：综合优化 & 决策变量、目标层级、硬约束、可行域和前后问接口；多目标的优先级或权重来源。 & 编码、初始解、越界/不可行处理、搜索范围、停止条件、重复协议和同口径基线。 & 目标值、约束残差、候选位置、成本与现实含义；有限改善要同时报最大值和平均值。\\
C：数据决策 & 原字段--清洗/聚合--特征--标签--划分--输出谱系；指标方向、业务语义和缺失类型。 & 外层测试冻结，预处理只读训练折；模型损失、阈值、类别不平衡、随机重复及预测到决策的接口。 & 训练/验证/测试角色、基线、混淆矩阵或误差、风险/收益和场景复评；同源指标不得冒充外部验证。\\
\bottomrule
\end{longtable}

\paragraph{A 类段落示例。}“在\slot{前期时段/稳定窗口}内，\slot{观测量}变化趋于\slot{形态}，因而先用\slot{方程/离散格式}描述\slot{状态}。当\slot{边界/事件}被触发时，计算转入\slot{局部求解}；将结果回代到\slot{守恒/几何条件}后，\slot{误差或边界结论}仍在\slot{范围}内。”

\paragraph{B 类段落示例。}“把\slot{资源/对象}直接平均分配会使\slot{具体约束/时段}出现\slot{违约或浪费}。因此保留\slot{硬约束}，先生成\slot{可行候选}，再用\slot{目标层级/求解器}比较。候选在同一\slot{历史或场景基准}下复评，最终选择\slot{方案}，代价是\slot{明确的牺牲}。”

\paragraph{C 类段落示例。}“原始字段\slot{名称}中，\slot{缺失/异常/重复}并不表示同一件事，故先按\slot{规则}分流，再在\slot{训练划分}内计算\slot{特征/尺度}。预测模型输出\slot{标签/概率}后，只有在\slot{阈值和业务规则}下才进入决策；在\slot{共同环境/独立测试}复评后，结论限定为\slot{范围}。”

\subsection{图、表与代码的可替换组件}

以下组件借鉴根目录培训版 \texttt{main.tex} 的排版做法\cite{friendtemplate}，但把示例数字和外部图片改成空位，避免把示例误当作本队结果。图表必须在正文中被提出、读数和解释；完整代码放附录或支撑材料，正文只留能说明算法接口的短片段。

\begin{figure}[htbp]
  \centering
  \fbox{\rule{0pt}{3.0cm}\rule{0.82\linewidth}{0pt}}
  \caption{\slot{图的对象、条件或结果名称}}
  \label{fig:template-placeholder}
\end{figure}

\begin{table}[htbp]
  \centering
  \caption{\slot{结果表名称}}
  \label{tab:template-result}
  \begin{tabular}{p{3.2cm}p{3.0cm}p{3.0cm}p{3.2cm}}
    \toprule
    对象/方案 & 指标一（单位） & 指标二（单位） & 结论所需状态\\
    \midrule
    \slot{A} & \slot{数值} & \slot{数值} & \slot{可行/超限/排名}\\
    \slot{B} & \slot{数值} & \slot{数值} & \slot{……}\\
    \bottomrule
  \end{tabular}
\end{table}

\begin{lstlisting}[language=Python,caption={分问代码入口与输出接口示例},label={lst:template-code}]
# 1. read_data: read the same files listed in the appendix
# 2. preprocess: keep the rule, unit and split used in the manuscript
# 3. solve_problem: return solution, objective and constraint residuals
# 4. validate: run the predeclared test and save the table/figure
data = read_data("[input file]")
state = preprocess(data, config)
result = solve_problem(state, config)
validate(result, state, config)
\end{lstlisting}

代码说明不要只写函数名。把\slot{输入对象}、\slot{返回对象}、\slot{随机种子/参数}和\slot{输出文件}填入附录复现清单；若某一步失败，保留状态和原因，而不是在导出结果前静默改动。

\subsection{结果分析与模型检验：表后必须继续说话}

\begin{longtable}{p{2.2cm}p{3.0cm}p{3.0cm}p{3.0cm}p{3.5cm}}
\caption{结果解释台账}\label{tab:template-result-ledger}\\
\toprule
结果对象 & 条件/对象 & 读数与单位 & 与谁比较 & 原因、影响和边界\\
\midrule
\endfirsthead
\toprule
结果对象 & 条件/对象 & 读数与单位 & 与谁比较 & 原因、影响和边界\\
\midrule
\endhead
\slot{表/图/输出编号} & \slot{情景、样本、时段} & \slot{数值、方向、数量级} & \slot{基线/候选/实测} & \slot{机制、约束、现实动作和不可外推范围}\\
\slot{} & \slot{} & \slot{} & \slot{} & \slot{}\\
\bottomrule
\end{longtable}

每张表或图后至少补一段：“在\slot{条件}下，\slot{对象}的\slot{指标}为\slot{数值}，相对于\slot{基准}变化\slot{幅度}。变化主要出现在\slot{局部对象/时段}，与\slot{方程项、数据特征、约束状态或决策规则}相符，因此对\slot{现实任务}的含义是\slot{动作}；在\slot{边界}之外不作同样外推。”

\begin{longtable}{p{2.6cm}p{3.4cm}p{3.0cm}p{3.0cm}p{2.8cm}}
\caption{模型检验台账}\label{tab:template-validation-ledger}\\
\toprule
检验问题 & 独立证据/数据 & 指标与接受标准 & 结果 & 裁决及影响范围\\
\midrule
\endfirsthead
\toprule
检验问题 & 独立证据/数据 & 指标与接受标准 & 结果 & 裁决及影响范围\\
\midrule
\endhead
\slot{守恒/回代是否成立} & \slot{未参与拟合的量/公式} & \slot{误差定义、阈值} & \slot{数值} & \slot{保留/修正哪一层}\\
\slot{预测/分类是否可用} & \slot{独立测试或交叉验证} & \slot{MAE/RMSE/召回等} & \slot{} & \slot{}\\
\slot{方案是否可执行} & \slot{约束复算/场景} & \slot{违约率、成本、风险} & \slot{} & \slot{}\\
\bottomrule
\end{longtable}

检验实施前先写判据：“若\slot{观测结果/指标}满足\slot{阈值}，则保留\slot{模型或参数}；否则只修改\slot{对应关系}并重新计算。”训练回代、内层调参和最终独立测试要分开；一次重复运行不能写成“全局最优”，一条相关曲线不能写成因果验证。

\subsection{灵敏度与稳健性：改变什么、为什么改变、何时换方案}

\begin{longtable}{p{2.4cm}p{3.0cm}p{3.0cm}p{3.0cm}p{3.5cm}}
\caption{灵敏度/稳健性台账}\label{tab:template-sensitivity-ledger}\\
\toprule
扰动量 & 基准与来源 & 扰动范围/固定量 & 输出变化 & 决策切换与结论边界\\
\midrule
\endfirsthead
\toprule
扰动量 & 基准与来源 & 扰动范围/固定量 & 输出变化 & 决策切换与结论边界\\
\midrule
\endhead
\slot{参数/误差/场景} & \slot{测量、估计或题面来源} & \slot{范围、步长、其余固定} & \slot{最大/平均/排序/曲线} & \slot{是否换方案，阈值在哪里}\\
\slot{} & \slot{} & \slot{} & \slot{} & \slot{}\\
\bottomrule
\end{longtable}

正文句式：“以\slot{基准}为中心，只改变\slot{参数}，其余量保持\slot{固定规则}。当\slot{参数}位于\slot{范围}时，\slot{输出}变化\slot{幅度/方向}；超过\slot{阈值}后，\slot{约束/排序/风险}发生切换。因此本文把\slot{方案或参数}限定在\slot{范围}，而不把一次基准运行推广到全部情景。”

\subsection{模型评价与改进：优缺点必须落到对象}

\begin{longtable}{p{2.3cm}p{4.0cm}p{3.6cm}p{3.8cm}}
\caption{模型评价与改进台账}\label{tab:template-evaluation-ledger}\\
\toprule
评价方面 & 已有证据 & 具体优点或不足 & 可执行的改进及新检验\\
\midrule
\endfirsthead
\toprule
评价方面 & 已有证据 & 具体优点或不足 & 可执行的改进及新检验\\
\midrule
\endhead
精度/解释 & \slot{误差、回代或图表} & \slot{对哪些对象有效/在哪里失效} & \slot{新增数据、关系或校准方法}\\
计算/复现 & \slot{复杂度、运行记录、代码} & \slot{瓶颈和依赖} & \slot{缓存、分层求解或更换离散格式}\\
现实性/外推 & \slot{约束和场景复评} & \slot{未观测对象、时窗或风险} & \slot{追加样本、边界场景和接受阈值}\\
\bottomrule
\end{longtable}

不要把“速度快、精度高、适用范围广”单独列作优点。改进项必须写成“新增\slot{数据/关系}，修改\slot{模型层}，预期检查\slot{指标/边界}”；现有数据没有支撑的改进只能标为后续工作，不能倒写成本文已经完成。

\subsection{参考文献与附录：让读者能够追溯}

正文首次使用外部定义、算法、数据或软件时引用对应来源；文后不放未使用的条目。附录按“输入--处理--求解--输出”排列，并至少填下面的清单：

\begin{longtable}{p{3.0cm}p{4.2cm}p{4.2cm}p{3.2cm}}
\caption{附录复现清单}\label{tab:template-appendix-ledger}\\
\toprule
项目 & 必须给出的内容 & 与正文的对应关系 & 验收状态\\
\midrule
\endfirsthead
\toprule
项目 & 必须给出的内容 & 与正文的对应关系 & 验收状态\\
\midrule
\endhead
数据字典 & 文件名、字段、单位、缺失/异常编码、主键 & 问题重述、预处理和符号表 & \slot{待核/已核}\\
代码入口 & 主脚本、函数调用顺序、依赖和版本 & 模型求解段、输出表/图 & \slot{}\\
参数与随机性 & 初值、范围、步长、种子、重复次数、停止条件 & 模型卡和灵敏度表 & \slot{}\\
中间结果 & 关键状态、候选、约束残差、日志或全量表位置 & 结果分析和检验 & \slot{}\\
最终文件 & 运行命令、输出路径、文件哈希或版本日期 & 摘要数字和交稿文件 & \slot{}\\
\bottomrule
\end{longtable}

\subsection{正式提交前的十四项验收单}

\begin{enumerate}
  \item 题面对象、单位、时间窗和硬约束与附件逐项对上；
  \item 每个分问都有输入、输出和交给下一问的接口；
  \item 问题分析中至少保留一个由数据、边界、失败试算或物理关系触发的实际取舍；
  \item 每条模型假设都说明计算后果和失效范围；
  \item 符号、变量域、下标、单位与公式和代码一致；
  \item 目标/方程、约束、边界、初值和求解器可以互相追溯；
  \item 算法段写明初始化、更新、不可行处理、停止条件和输出；
  \item 结果表图不是装饰，正文逐项解释读数、原因、影响和边界；
  \item 检验独立于拟合或寻优证据，并写明指标、阈值和裁决；
  \item 灵敏度固定共同基准，写清扰动来源和决策切换；
  \item 预测值、评级或综合分进入决策时，误差传播和接口没有被省略；
  \item A/B/C 题型专用插槽已按实际题目填入，不把算法名堆成清单；
  \item 参考文献逐条有正文出处；AI 工具使用声明位于参考文献之前，使用过 AI 时支撑材料含\texttt{AI 工具使用详情.pdf}；
  \item 附录运行入口、参数、依赖和输出齐全；
  \item 删除所有\slot{填空位}、训练提示和未核对数字后，再做一次全文、代码和 PDF 的同源复算。
\end{enumerate}

\appendix
\section{模板提交前速查}

\begin{enumerate}
  \item 摘要中的结论、正文表格和附件输出使用同一数值与单位；
  \item 每个模型都有变量、方程或目标、约束、求解入口和结果解释；
  \item 每幅图和每张表在正文中被提出，并至少承担一次读数或比较；
  \item 关键结论配有回代、误差、交叉验证、基线、反例或物理边界中的一种检查；
  \item 附录列出运行顺序、依赖、随机种子、输入文件和输出文件。
\end{enumerate}

\section*{材料来源说明}
\addcontentsline{toc}{section}{材料来源说明}

\begin{thebibliography}{9}
\bibitem{format2026}
全国大学生数学建模竞赛组委会：《全国大学生数学建模竞赛论文格式规范（2026年修订稿）》，
工作区文件\texttt{全国大学生数学建模竞赛论文格式规范（2026年修订稿）.pdf}，2026 年 8 月 11 日核读；
\url{https://www.mcm.edu.cn/html_cn/node/4cd596519c9eb9fbd866398f6df0caa3.html}。

\bibitem{rules2026}
全国大学生数学建模竞赛组委会：《全国大学生数学建模竞赛参赛规则（2026年修订稿）》，
工作区文件\texttt{全国大学生数学建模竞赛参赛规则（2026年修订稿）.pdf}，2026 年 8 月 11 日核读；
\url{https://www.mcm.edu.cn/html_cn/node/9d8e511fe7a1447b35f53a82c908e2e0.html}。

\bibitem{ai2026}
全国大学生数学建模竞赛组委会：《全国大学生数学建模竞赛人工智能工具使用规定（2026年试行）》，
工作区文件\texttt{全国大学生数学建模竞赛人工智能工具使用规定 （2026年试行）.pdf}，
2026 年 8 月 11 日核读。

\bibitem{corpus59}
2020--2025 年全国大学生数学建模竞赛 A/B/C 题编号论文语料，59 篇、2892 页；
写作 Skill 的证据索引 \texttt{references/corpus-index.md}，
结构统计版本 2026 年 8 月 1 日。

\bibitem{friendtemplate}
数学建模竞赛论文写作培训模板，工作区根目录文件 \texttt{main.tex}，2026 年 8 月 10 日核读。

\bibitem{mindmap}
数学建模学习思维导图，工作区文件 \texttt{1.png}，2026年7月核读。

\bibitem{A053}
\emph{基于目标优化的板凳龙形态调节}，2024年全国大学生数学建模竞赛A题一等奖论文语料，
工作区文件 \texttt{2024/A053.pdf}。

\bibitem{problem2024A}
2024年全国大学生数学建模竞赛A题：板凳龙，工作区赛题原文。
\end{thebibliography}

\end{document}
